%%%%%%%%%%%%%%%%%%%%%%%%%%%%%%%%%%%%%%%%%
% University Assignment Title Page 
% LaTeX Template
% Version 1.0 (27/12/12)
%
% This template has been downloaded from:
% http://www.LaTeXTemplates.com
%
% Original author:
% WikiBooks (http://en.wikibooks.org/wiki/LaTeX/Title_Creation)
%
% License:
% CC BY-NC-SA 3.0 (http://creativecommons.org/licenses/by-nc-sa/3.0/)
% 
% Instructions for using this template:
% This title page is capable of being compiled as is. This is not useful for 
% including it in another document. To do this, you have two options: 
%
% 1) Copy/paste everything between \begin{document} and \end{document} 
% starting at \begin{titlepage} and paste this into another LaTeX file where you 
% want your title page.
% OR
% 2) Remove everything outside the \begin{titlepage} and \end{titlepage} and 
% move this file to the same directory as the LaTeX file you wish to add it to. 
% Then add \input{./title_page_1.tex} to your LaTeX file where you want your
% title page.
%
%%%%%%%%%%%%%%%%%%%%%%%%%%%%%%%%%%%%%%%%%
%\title{Title page with logo}
%----------------------------------------------------------------------------------------
%	PACKAGES AND OTHER DOCUMENT CONFIGURATIONS
%----------------------------------------------------------------------------------------

\documentclass[12pt]{article}

\usepackage[francais]{babel}
\usepackage[utf8x]{inputenc}
\usepackage[T1]{fontenc}
\usepackage{color}

\usepackage{amsmath}
\usepackage{graphicx}
\usepackage{enumerate}
\usepackage{url}
\usepackage{listings}

\usepackage{xcolor}
\colorlet{keyword}{blue!100!black!80}
\colorlet{comment}{green!90!black!90}
\lstdefinestyle{vhdl}{
   language     = VHDL,
   basicstyle   = \ttfamily,
   keywordstyle = \color{keyword}\bfseries,
   commentstyle = \color{comment}
}

\lstdefinestyle{equ}{
   basicstyle   = \ttfamily,
   keywordstyle = \color{keyword}\bfseries,
   commentstyle = \color{comment}
}

% Define new command
\newcommand{\HRule}{\rule{\linewidth}{0.5mm}}

\newcommand{\crt}{\emph{Nexys 4 DDR\ }}
%\def\thesubsection{\alph{section}}

\begin{document}

\begin{titlepage}

\center % Center everything on the page
 
%----------------------------------------------------------------------------------------
%	HEADING SECTIONS
%----------------------------------------------------------------------------------------

\textsc{\LARGE Universit\'e Pierre et Marie Curie}\\[1.5cm] % Name of your university/college
\textsc{\Large PSESI}\\[0.5cm] % Major heading such as course name

%----------------------------------------------------------------------------------------
%	TITLE SECTION
%----------------------------------------------------------------------------------------

\HRule \\[0.4cm]
{ \huge \bfseries Projet Centrale DCC}\\[0.4cm] % Title of your document
{ \huge \bfseries Document Définition Interface et Logique Enclenchement}\\[0.4cm] % Title of your document
\HRule \\[1.5cm]
 
%----------------------------------------------------------------------------------------
%	AUTHOR SECTION
%----------------------------------------------------------------------------------------

\begin{minipage}{0.4\textwidth}
\begin{flushleft} \large
\emph{\'Etudiant:}\\
Maxime \textsc{AYRAULT} 3203694 % Your name
\end{flushleft}
\end{minipage}
~
\begin{minipage}{0.4\textwidth}
\begin{flushright} \large
\emph{Encadrant:} \\
Julien \textsc{DENOULET} % Supervisor's Name
\end{flushright}
\end{minipage}\\[2cm]

%----------------------------------------------------------------------------------------
%	DATE SECTION
%----------------------------------------------------------------------------------------

{\large \today}\\[2cm] % Date, change the \today to a set date if you want to be precise

%----------------------------------------------------------------------------------------
%	LOGO SECTION
%----------------------------------------------------------------------------------------

%%\begin{figure}
%%  \subfigure[]{\includegraphics[scale=0.2\textwidth]{logo.png}} 
%%\end{figure} 
\includegraphics[width=0.2\textwidth]{logo.png}

%----------------------------------------------------------------------------------------

\vfill % Fill the rest of the page with whitespace

\end{titlepage}



%\begin{abstract}
%Your abstract here.
%\end{abstract}

\section{Introduction}
\label{sec:introduction}
Présentation des 3 interfaces~:
\begin{itemize}
\item Interface Homme Machine
\item Interface Centrale DCC et le circcuit ferroviaire
\item Interface interne Centrale DCC - Moteur IXL
\end{itemize}

et présentation des règles de la logique d'enclenchement.


\section{Interface Homme Machine}
\label{sec:int-dcc}
Sur la carte
\begin{itemize}
\item Bouton gauche -> changement du paramètre
\item Bouton droit  -> changement de la valeur du paramètre
\end{itemize}

\hrulefill

Dans le design
\begin{enumerate}
\item en sortie :
  \begin{itemize}  
  \item les 7 selecteurs de segment Ca -> Cg
  \item les 8 selecteurs de digit   An (7 downto 0)
  \item octet adresse \textbf{\emph{(addr or aigui)}}  dans ADD
  \item octet speed \textbf{\emph{(speed)}}  dans SPEED
  \item octet fonction \textbf{\emph{(feat)}}  dans FCTN    
  \end{itemize}  
\item en entrée :
  \begin{itemize}  
  \item une horloge à 100 MHz
  \item un reset synchrone actif haut 
  \item la valeur du bouton gauche BTNL
  \item la valeur du bouton droit BTNR    
  \end{itemize}  
\end{enumerate}

\subsection{\underline{Interface commande train}}
\label{sec:ihm-train}

\bigskip

8 vitesses
\begin {enumerate}
\item vitesse 0 -> \emph{``01100000''}
\item vitesse 1 -> \emph{``01100010''}
\item vitesse 2 -> \emph{``01100100''}
\item vitesse 3 -> \emph{``01101000''}
\item vitesse 4 -> \emph{``01101010''}
\item vitesse 5 -> \emph{``01110000''}
\item vitesse 6 -> \emph{``01110100''}
\item vitesse 7 -> \emph{``01111111''}
\end {enumerate}  

\bigskip

6 adresses de train
\begin {enumerate}
\item adresse 0 -> \emph{``00000000''}  
\item adresse 1 -> \emph{``00000001''}
\item adresse 2 -> \emph{``00000010''}
\item adresse 3 -> \emph{``00000011''}
\item adresse 4 -> \emph{``00000100''}
\item adresse 5 -> \emph{``00000101''}
\item adresse 6 -> \emph{``00000110''}
\end {enumerate}  

\bigskip

8 adresses de d'aiguillage
\begin {enumerate}
\item elles ne sont pas encore définies. mais de 10 à 18 pour faire simple.
\end {enumerate}  

\bigskip

X fonctions
\begin {enumerate}
\item phares on     -> \emph{``10010000''}
\item bruit moteur  -> \emph{``10000001''}
\item klaxon 1 long -> \emph{``10000100''}
\item klaxon 1 bref -> \emph{``10100100''}
\item bruit du vent -> \emph{``10100001''}
\end {enumerate}  

\bigskip

\subsection{\underline{Interface itinéraire}}
\label{sec:ihm_iti}

bla bla bla...

\newpage

\section{Interface Centrale DCC - Train, aiguillage et capteurs}
\label{sec:int_dcc}

\subsection{\underline{Interface Train}}
\label{sec:int_train}

\bigskip
Trames envoyéees 4 par 4 séparées chacunes par 400 us avec valeur 0.
\begin{center}
\begin{tabular}{|l|l|l|}
  \hline
  1er trame & IDLE \\
  \hline
  2eme trame & Vitesse \\
  \hline
  3eme trame & Action \\
  \hline
  4eme trame & IDLE \\
  \hline
\end{tabular}
\end{center}

\medskip
Une fois une trame envoyée, la répéter toutes les 55 ms. Si aucune
trame de chargée alors envoyer des trames idle.

\medskip
\begin{center}
\begin{tabular}{lll}
adresse & => & ``11111111'' \\
donnée  & => & ``00000000'' \\
CRC     & => & ``11111111'' \\
\end{tabular}
\end{center}

\medskip
Paramètres à envoyer stockés dans \emph{Master}.


-> à chaques pression sur le bouton central changement de la valeur
des paramètres à envoyer, Les paramètre viennent de l'IHM et
remplacent la valeur IDLE et sont envoyés en continue tant que pas de
nouvelle prssion sur bouton C.

\bigskip
\subsection{\underline{Interface Aiguillage}}
\label{sec:int_aig}

bla bla bla

\subsection{\underline{Interface Capteurs}}
\label{sec:int_cap}

bla bla bla...


\newpage
\section{\underline{Interface Centrale DCC - Moteur IXL}}
\label{sec:int_ixl}   

La logique d'enclenchement est basée sur les notions de signaux, de
circuit de voie et d'itinéraire ferroviaire.

Un \emph{signal} est le moyen de communication entre le conducteur du
train et la logique d'enclenchement.  Il permet d'indiquer au
conducteur l'état de la voie se trouvant devant le train (par
définition un train ne recule jamais, il est nécessaire de changer de
cabine de conduite pour changer la direction du train). Il existe
beaucoup de signaux ferroviaire. Les trois plus importants sont~:
\begin{itemize}
\item signal \emph{rouge} qui signifie que le train doit s'arrêter
  devant le signal car la voie devant lui n'est pas sûre (présence
  d'un autre train, aiguillage en train d'être commandé, panne
  matériel...)
\item signal \emph{jaune} qui signifie qu'il y a un danger à
  proximité. Le train peut passer le signal mais à vitesse réduite
  seulement et doit être capable de s'arrêter sur une courte
  distance. Ceci permet de fluidifier le traffic.
\item signal \emph{vert} qui signifie que la voie devant le train est
  libre de tout danger et que le train peut poursuivre sa route en
  respectant les vitesses limites.
\end{itemize}

L'objectif des opérateurs consiste à toujours présenter un signal vert
devant un train de façon à ce que le train n'est jamais à s'arrêter à
cause d'un danger.


\medskip
Le projet implémente le signal rouge et le signal vert. Le signal
jaune ne peut pas être implémenté sans conducteur dans le train. Il
n'y a pas de signaux sur le circuit ferroviaire, mais la notion de
signaux sera quand même implémentée. Un signal \emph{virtuel} sera
créé pour chaque circuit de voie. Il autorisera ou pas l'entrée d'un
train dans le circuit de voie auquel il est associé.  Ceci permettra
de commander les trains. 

\medskip
Un \emph{Ciruit de Voie} (CdV) représente la plus petite partie de
voie sur lequel il est possible de détecter la présence d'un
train. Pour un train réel, la détection de la présence d'un train
s'effectue par les essieux du train. En effet, les essieux du train
sont conducteur (métal), ils peuvent servir comme un
\emph{interupteur} entre les 2 rails de la voie. Si une partie du
train est présent alors les 2 rails sont au même potentiel (le circuit
est dit fermé), sinon les rails sont a des potentiels différents (le
circuit est dit ouvert).

\medskip
Pour le projet, un Circuit de voie sera délimité par 2 capteurs de
passage(sauf en zone d'aiguille où le circuit de voie sera délimité
par les capteurs de chaque branche). L'état d'occupation du CdV sera
fourni pas le sens de parcours identifié par le capteur pour chaque
train.

\medskip
Un \emph{itinéraire ferroviaire} représente une portion de voie entre
2 signaux. Il contient les mouvements d’aiguilles nécessaires pour que
le train parcourt le chemin entre le signal de départ et le signal de
sortie de l'itinéraire. Dans le système d'enclenchement, un itinéraire
est représenté par une machine à états. Les états possibles sont:
\begin{itemize}
\item \emph{commandé}. L’opérateur demande au train de parcourir
  l’itinéraire choisi. Il a par conséquent demandé la formation de
  l’itinéraire au système d’enclenchement. Si les conditions sur les
  aiguilles à manoeuvrer pour former l’itinéraire sont sûres, le
  système d’enclenchement passe à l’état de formation de l’itinéraire
  sinon il est rejeté.
\item \emph{formé}. Les aiguilles sont positionnées dans la bonne
  direction. Il n’est plus possible de bouger les aiguilles tant que
  le train se situe dans une certaine zone.
\item \emph{détruit}. Le train a parcouru l’itinéraire (il a passé le
  signal de sortie), l'itinéraire peut être détruit. Ou l'opérateur a
  demandé la destruction de l'itinéraire (en cas d'erreur par exemple)
  Dans ce cas, si le train ne se trouve pas trop proche du signal
  d'entrée alors l'itinéraire peut être détruit sinon la demande de
  destruction est rejetée.  Ceci permet d'autoriser le parcours sur
  les aiguilles et la zone d’un autre train.
\end{itemize}


L'interface entre la Centrale DCC et le Moteur IXL est bi-directionnelle. 
Dans le sens Centrale vers IXL, les informations envoyées sont~:
\begin{itemize}
\item les commandes d'itinéraire
\item l'information des capteurs de passage
\item l'état des aiguilles (gauche/droite)
\end{itemize}

Dans le sens IXL vers Centrale, les informations envoyées sont~:
\begin{itemize}
\item les commandes des aiguilles
\item l'état des itinéraires (commandé/formé/détruit)
\item l'état des signaux (rouge/vert)
\end{itemize}

\subsection{\underline{Commande itinéraire}}
\label{sec:ixl_iti}

bla bla bla...

\subsection{\underline{\'Information capteur de passage}}
\label{sec:ixl_cdv}


\medskip
La déclaration VHDL
\begin{lstlisting}[style=vhdl]
  type train_addr is 7 downto 0 ;
  type train_dir is bit ;

  type sensor is
   record
       addr : train_addr;
       dir  : train_dir;
   end record;

  type SE_state is array (31 downto 0) of sensor;
\end{lstlisting}


\subsection{\underline{\'Etat des aiguillages}}
\label{sec:ixl_aig}

L'état d'une agiuille est déterminé par la combinaison entre les 2
capteurs d'aiguilles.  Une aiguille peut avoir 3 états:
\begin{itemize}
\item \emph{gauche} lorsque le contact gauche est à l'état 1
\item \emph{droite} lorsque le contact droit est à l'état 1
\item \emph{en mouvement} lorsqu'aucun des 2 contacts ne sont à 1
\end{itemize}

Il n'y a pas de capteur fournissant l'état d'une aiguille sur la
maquette.  Par conséquent, l'état d'une aiguille correspondra à l'état
de la dernière commande envoyée à l'aiguille. L'état \emph{en
  mouvement} ne sera pas considéré.
 
L'état d'une aiguille sera représenté par 1 bit avec la convention
\begin{itemize}
\item 0 = droite
\item 1 = gauche
\end{itemize}
 
Il y a 4 aiguilles sur le circuit. L'indice du tableau correspond au
numéro de l'aiguille.

\medskip
La déclaration VHDL
\begin{lstlisting}[style=vhdl]
  type Sw_state is array (7 downto 0) of bit;
\end{lstlisting}


\subsection{\underline{\'Etat des itinéraires}}
\label{sec:st_iti}

Les états possibles pour un itinéraire sont; commandé, formé, détruit
et rejeté. Il sera représenté par 2 bits avec la conversion
\begin{itemize} 
\item 00 = commandé
\item 01 = formé
\item 10 = détruit
\item 11 = rejeté
\end{itemize}

Il y a 32 itinéraires possibles. L'indice dans le tableau correspond
au numéro d'itinéraire.

\medskip
La déclaration VHDL
\begin{lstlisting}[style=vhdl]
  type Ro_state_t is 3 downto 0 ;

  type Ro_state is array (31 downto 0) of Ro_state_t;
\end{lstlisting}


\subsection{\underline{Commande des aiguillages}}
\label{sec:cmd_aig}

bla bla bla...

\subsection{\underline{\'Etat des signaux}}
\label{sec:st_sig}

bla bla bla...

  
\subsection{\underline{\'Etat des Circuits de Voie}}
\label{sec:st_sig}

Chaque circuit de voie possède 2 états; \emph{Libre} ou
\emph{Occupé}. De façon à pouvoir effectuer la gestion des trains, il
est nécessaire d'identifier quel train occupe le circuit de voie.  Il
y a 6 trains au maximum. Les états sont codés come suit~:
\begin{itemize}
  \item libre : 0x07
  \item occupé : numéro du train 0x01 à 0x06
\end{itemize}  

Il y a 20 circuits de voie sur la maquette de train. L'indice du
tableau correspond au numéro du Circuit de Voie.




\newpage
\section{\underline{Logique Enclenchement}}
\label{sec:log_enc}

Cette secion présente la logique d'enclenchement, sa justification et
les conversions de nommage utilisées dans les équations.

\subsection{\underline{Circuit de Voie}}
\label{sec:CdV}

Un Circuit de Voie est la section de voie se trouvant entre 2 capteurs
de présence consécutifs. Un capteur de présence permettent de
connaître, le numéro du train et son sens de déplacement. Le sens de
déplacement va permettre de savoir quel Circuit de voie occupé par le
train. Deux directions sont arbitrairement déclarées pour chaque
capteur; Up et Down.

Les convention de nommage des éléments~:
\begin{itemize}
\item $SE\_Up_{nn}$, 1 lorsqu'un train vient de passer le capteur numéro
  nn dans le sens Up
\item $SE\_Do_{nn}$, 1 lorsqu'un train vient de passer le capteur numéro
  nn dans le sens Down
\item $TC_{nn}$, 1 lorsque le Circuit de Voie est libre (aucun train
  présent sur le CdV), 0 lorsque le circuit de voie est occupée. Cette
  logique va dans le sens de la sécurité, en cas de perte de courant
  ou de dysfonctionnement, il n'est pas possible de savoir si le
  Circuit de VOie est occupé ou libre, donc on choisit la valeur la
  plus sécuritaire; Occupé même si aucun train n'est présent.
\end{itemize}

\'A l'initialisation, tous les Circuits de Voie sont libres.

\medskip
Le Circuit de voie $TC_{nn}$ est occupé si une des condition suivantes
est vérifiée~:
\begin{itemize}
\item le capteur $SE\_Up_{nn}$ passe à 1, un train vient de rentrer dans
  le sens Up
\item le capteur $SE\_Do_{mm}$ passe à 1, un train vient de rentrer dans
  le sens Down
\item le Circuit de Voie est occupé et le train n'est pas sorti dans
  le sens Up ou Down
\end{itemize}

Il faut inverser cette logique car lorsque le Circuit de voie est
occupé, on cherche à avoir une valeur 0.

\medskip
L'équation générale d'occupation d'un Circuit de Voie est~:
$$\boxed{
  TC_{nn} \Leftarrow \neg SE\_Up_{nn} \land \neg SE\_Do_{mm} \land
  (TC_{nn} \lor SE\_Do_{nn} \lor SE\_Up_{mm});
}$$


\subsection{\underline{Signaux}}
\label{sec:esp}

Il existe deux types de signaux; les signaux d'espacement qui
garantissent que 2 trains ne vont pas se percuter et les sinaux de
manoeuvre qui garantissent qu'un train ne va pas parcourir une
aiguille qui n'est pas contrôlée.

Les convention de nommage des éléments~:
\begin{itemize}
\item $SI_{nn}$, 1 lorsque le signal est Vert, 0 pour afficher un signal
  Rouge
\end{itemize}

\'A l'initialisation, tous les signaux sont au Vert.

Un Signal d'espacement passe~:
\begin{itemize}
\item au Rouge, lorsque le signal d'espacement ou le signal de
  manoeuvre est au Rouge
\item au Vert, lorsque le signal d'espacement et le signal de
  manoeuvre sont au Vert
\end{itemize}


\medskip
L'équation générale d'un signal est~:
$$\boxed{
  SI_{nn} \Leftarrow SI\_SP_{nn} \land SI\_IX_{nn};
}$$


\subsection{\underline{Signal d'espacement des trains}}
\label{sec:esp}

Il est dangereux d'avoir 2 trains dans le même Circuit de Voie car
rien ne permet de garantir que les 2 trains ne vont pas se
percuter. Il n'est pas certain qu'un conducteur puisse arrêter le
train devant un signal car la distance de freinage entre le signal et
le Circuit de Voie à protéger n'est pas toujours suffisante. C'est
pourquoi, il est nécessaire de toujours avoir un Circuit de Voie libre
entre 2 trains. Ce Circuit de Voie est nommé, \emph{CdV tampon}.  Dans
le cas de 2 trains qui se suivent, il est nécessaire que la longueur
du Circuit de Voie soit plus grande que la distance d'arrêt du train
suiveur.  Dans le cas de 2 trains en sens inverse, il est nécessaire
de s'assurer que la longueur du Circuit de Voie tampon est égale à 2
fois la longueur d'arrêt d'un train.

L'entrée dans un Circuit de Voie est protégée par un signal dit
\emph{d'espacement}.

Les convention de nommage des éléments~:
\begin{itemize}
\item $SI\_SP_{nn}$, 1 lorsque le signal est Vert, 0 pour afficher un
  signal Rouge
\end{itemize}

\'A l'initialisation, tous les signaux sont au Vert.

Un Signal d'espacement passe~:
\begin{itemize}
\item au Rouge, lorsque le $TC_{nn}$ est occupé ou le $TC_{mm}$ est occupé
\item au Vert, lorsque les $TC_{nn}$ et $TC_{mm}$ sont libres.
\end{itemize}

\medskip
L'équation générale d'un signal d'espacement est~:
$$\boxed{
  SI\_SP_{nn} \Leftarrow TC_{nn} \land TC_{mm};
}$$


\subsection{\underline{Signal de manoeuvre}}
\label{sec:esp}

Un signal de manoeuvre garantit qu'un train peut traverser une zone
d'aiguillage sans danger. C'est à dire que les aiguilles de la zone
sont toutes enclenchées.




\subsection{\underline{Zone d'aiguille}}
\label{sec:aig}

Une aiguille peut être commandée si elle n'est pas enclenchée et si il
n'y a pas de train sur la zone.


\subsection{\underline{Itinéraire}}
\label{sec:iti}

Bla bla bla...

\subsection{\underline{zone d'aiguille}}
\label{sec:aig}

Bla bla bla...



\newpage

\bibliographystyle{plain}
\bibliography{biblio}

\end{document}
